%% ch04_method.tex  —  Chapter 4: The SgDP Framework and HSP-Agile Instantiation (Step 1 polish staging)
%% HSP-Agile Thesis, East China Normal University, 2026
%% Staging rewrite: paper/hsp-agile/_rewrite/ch04_method.tex

\chapter{The \textsc{SgDP} Framework and HSP-Agile Instantiation}
\label{ch:method}

This chapter realises the two mechanisms introduced in
Chapter~\ref{ch:intro}: structured scenario-aware repair feedback~(C2) and
dual-gate acceptance~(C3), built on the SpecIR engineering hinge.
Chapter~\ref{ch:formalization} fixed the formal objects; here they become a
loop that an LLM can participate in.
Adapters and FSF lowering are necessary engineering for the HSP-Agile
instantiation and are not claimed as contributions.

The walkthrough opens with the \emph{primary} vignette from
Chapter~\ref{ch:intro}---the catch-all / \texttt{OTHERWISE} tier classifier---so
readers see the four stages as artefact transforms before the interfaces are
named.
Ordered-overlap examples (\code{classify\_signal}, \code{compute\_tax}) serve as
a \emph{secondary} thread for witness priority formulas $\Phi_i$.
Notation is inherited:
$t = (\mathcal{S}_t, \mathcal{X}_t, \mathcal{Y}_t)$, oracle $g_t$, and
first-match semantics over $\mathcal{S}_t = \langle s_1,\ldots,s_n\rangle$.

\section{Pipeline Intuition}
\label{sec:method:intuition}

The pipeline synthesises from FSF text, checks on first-match SMT witnesses,
repairs from a typed failure record, and releases only when every witness
matches \emph{and} no high-severity structural pattern remains.
Figure~\ref{fig:sgdp-overview} (Chapter~\ref{ch:intro}) depicts that loop; the
dashed orange arc marks budget~$K=3$.

\section{Worked Example: Catch-All Rescue End to End}
\label{sec:method:worked}

\paragraph{Primary vignette.}
Listings~\ref{lst:others-fsf}--\ref{lst:others-bug} define a three-scenario
tier classifier: Reject on negative scores, Pass at or above a floor, and
\texttt{OTHERWISE} must return \texttt{Review}.
The buggy candidate collapses the residual into \texttt{Pass}.
Busy Reject/Pass tests all succeed, so the mid-band never reaches a human.

The walkthrough below shows how each stage transforms that failure into a
repairable record and, after one iteration, into $\mathrm{Accept}=1$.

\paragraph{Stage~1 (synthesis).}
The LLM receives the FSF text of Listing~\ref{lst:others-fsf} and emits a
Python function.
A common first attempt is Listing~\ref{lst:others-bug}: the named guards look
correct, but the catch-all returns the wrong constant.

\paragraph{Stage~2 (witnesses and check).}
The priority formulas for the three scenarios include a residual witness for
the catch-all:
\begin{align}
  \Phi_1 &\triangleq \mathit{score}<0, \\
  \Phi_2 &\triangleq \mathit{score}\ge\mathit{floor}
           \;\land\; \neg(\mathit{score}<0), \\
  \Phi_3 &\triangleq \neg(\mathit{score}<0)
           \;\land\; \neg(\mathit{score}\ge\mathit{floor}).
  \label{eq:phi-catchall}
\end{align}
Z3 returns, for $\Phi_3$, a mid-band model such as $(\mathit{score},\mathit{floor})=(3,10)$.
The oracle requires \texttt{"Review"}; the buggy candidate returns
\texttt{"Pass"}, so $\mathrm{Conf}(c_1,t)<1$.

\paragraph{Stage~3 (Semantic Feedback IR).}
The failure is recorded as a typed JSON object, then rendered to a repair
prompt.  In abbreviated form:
\begin{quote}
\textit{``For inputs $(\mathit{score}{=}3,\;\mathit{floor}{=}10)$,
your code returned \texttt{Pass}, but the specification requires
\texttt{Review}.
This input satisfies Scenario~S3 (\texttt{others}/\texttt{OTHERWISE}), the
catch-all residual after Reject and Pass.''}
\end{quote}

\paragraph{Stage~4 and acceptance.}
After repair restores the correct catch-all constant, every witness matches and
the pattern screen is clear, so $\mathrm{Accept}(c_2,t)=1$.
Figure~\ref{fig:accept-tree} makes the conjunctive decision explicit:
conformance alone does not release a candidate.

\begin{figure}[t]
  \centering
  \tikzinput{diagrams/tikz/accept_decision_tree}
  \caption{\textbf{Acceptance decision tree.}
    Release requires both full witness conformance and a clear high-severity
    pattern screen.
    Takeaway: high mean conformance does not imply acceptance when residual
    witnesses or patterns still fail.}
  \label{fig:accept-tree}
\end{figure}

\paragraph{Secondary vignette (ordering / $\Phi_i$).}
Ordered overlaps motivate the negated higher-priority guards inside
Equation~\eqref{eq:priority}.
On \code{classify\_signal}, $\Phi_1$ forces $\mathit{level}<0$ before the
Critical band; a buggy candidate that tests $\mathit{level}\ge\mathit{threshold}$
first returns \texttt{Critical} on $(-3,-10)$ where the oracle requires
\texttt{Error}.
Figure~\ref{fig:worked-example-trace} is that secondary artefact trail;
Listings~\ref{lst:tax-fsf}--\ref{lst:tax-bug} are the tax-bracket twin used in
Chapter~\ref{ch:intro}.

\begin{figure}[t]
  \centering
  \tikzinput{diagrams/tikz/worked_example_trace}
  \caption{\textbf{Secondary worked-example trail (\code{classify\_signal}).}
    Ordering bug on $\Phi_1$: FSF text yields witness $(-3,-10)$; the buggy
    candidate returns \texttt{Critical}; Semantic Feedback IR names the
    precedence violation; repair restores $S_1$-first evaluation.
    Primary vignette remains the catch-all tier classifier above.}
  \label{fig:worked-example-trace}
\end{figure}

For the secondary example the priority formulas are
\begin{align}
  \Phi_1 &\triangleq \mathit{level}<0, \label{eq:phi1-example}\\
  \Phi_2 &\triangleq \mathit{level}\ge\mathit{threshold}
           \;\land\; \neg(\mathit{level}<0), \label{eq:phi2-example}\\
  \Phi_3 &\triangleq \neg(\mathit{level}<0)
           \;\land\; \neg(\mathit{level}\ge\mathit{threshold}).
\end{align}

\section{The \textsc{SgDP} Framework}
\label{sec:method:sgdp}

\subsection{Framework Overview and Instantiation Interface}
\label{subsec:framework-overview}

The walkthrough above already used four ideas that are not FSF-specific:
an oracle derived from the specification, witnesses that hit first-match
regions, a structured failure record for repair, and a conjunctive release
gate.
Together they form the language-independent core of \textsc{SgDP}.

\begin{enumerate}
  \item \textbf{Specification Oracle} ($\mathcal{G}$).
    An executable oracle derived from the formal specification, without a
    reference implementation.
  \item \textbf{Constraint-guided Witness Generator} ($\mathcal{W}$).
    SMT-based witnesses aimed at the boundary and residual regions where
    conformance defects appear.
  \item \textbf{Scenario-aware Semantic Feedback IR} ($\mathcal{F}$).
    Structured failure records that carry specification-level semantics into
    repair prompts, instead of raw pass/fail bits.
  \item \textbf{Conjunctive Acceptance Predicate} ($\mathcal{A}$).
    A release gate that demands both oracle-grounded conformance and
    specification-derived quality criteria.
\end{enumerate}

Figure~\ref{fig:sgdp-overview} (Chapter~\ref{ch:intro}) remains the map.
An instantiation for language $\mathcal{L}$ provides:
(i)~Specification Adapter $\text{Adapt}_L : \text{Text}_L \to \text{SpecIR}$;
(ii)~oracle $g_t$ derived from $\text{SpecIR}(t).\text{cases}$ under first-match;
(iii)~witness generator $\mathcal{W}(t)$ from SpecIR guard formulas;
(iv)~optional lowerer $\text{Lower}_{\text{FSF}}$ (HSP-Agile) $\to$ FSF TaskSpec;
(v)~Semantic Feedback IR per \texttt{schemas/semantic\_feedback.schema.json},
    rendered by $\rho : \mathcal{F} \to \text{Prompt}$.
Synthesis, repair, the pluggable structural screen, and acceptance are
shared core components.

\begin{figure}[t]
  \centering
  \tikzinput{diagrams/tikz/sgdp_specir_architecture}
  \caption{\textbf{Key observation:} notation changes stop at the adapter;
    the shared pipeline never sees SOFL, Mini-Z, or Mini-StateMachine text.
    \textbf{How to read:} adapters into SpecIR; optional
    $\text{Lower}_{\text{FSF}}$ for FSF-shaped screening; full context in
    Figure~\ref{fig:sgdp-overview}.}
  \label{fig:specir-arch}
\end{figure}

\subsection{HSP-Agile: SOFL/FSF Instantiation}
\label{subsec:hsp-agile}

HSP-Agile instantiates \textsc{SgDP} over FSF~\cite{liu2004soflbook,li2022qrs_companion}.
Each scenario $s_i = (g_i, \phi_i)$ pairs a guard with a postcondition; under
\emph{first-match semantics} the output for $\mathbf{x}$ is $\phi_i(\mathbf{x})$ where
$i = \min\{j : g_j(\mathbf{x})\}$.
The FSF oracle requires no reference implementation:
\begin{equation}
  g_t(\mathbf{x}) = \phi_i(\mathbf{x}) \text{ where } i = \min\{j : g_j(\mathbf{x})\}.
\end{equation}

\subsection{Constraint-guided Witness Generation}
\label{subsec:witness-gen}

Random tests rarely land in thin first-match regions---neither catch-all
residuals nor multi-conjunct overlaps.
The witness generator needs one input per scenario that activates
that scenario \emph{and no higher-priority one}.
For scenario $s_i$, it solves the \emph{priority formula}
\begin{equation}
  \label{eq:priority}
  \Phi_i \;\triangleq\; \hat{g}_i(\mathbf{x}) \;\land\; \bigwedge_{j=1}^{i-1} \lnot\hat{g}_j(\mathbf{x}),
\end{equation}
where $\hat{g}_j$ is the Z3 encoding of $g_j$.
The negated higher guards are essential: without them, a model of $g_i$
alone may still belong to $s_j$ for some $j<i$, and a catch-all encoded as
\texttt{True} without negations would not isolate the residual.
The witness set is $\mathcal{D}(t) = \{\mathbf{x}_1^*, \ldots, \mathbf{x}_n^*\}$.

\begin{proposition}[Scenario-aware Witness Soundness]
\label{thm:witness-soundness}
Assume $\hat{g}_i$ soundly encodes $g_i$ and Z3 returns a model $\mathbf{x}^*$ of
$\Phi_i$.
Then $g_i(\mathbf{x}^*)$ holds and $g_j(\mathbf{x}^*)$ is false for all $j < i$;
hence $g_t(\mathbf{x}^*) = \phi_i(\mathbf{x}^*)$.
\end{proposition}
\begin{proof}
By the encoding assumption and Z3 soundness,
$g_i(\mathbf{x}^*)$ and $\lnot g_j(\mathbf{x}^*)$ for every $j<i$.
By Definition~\ref{def:oracle},
$i=\min\{j:g_j(\mathbf{x}^*)\}$, so
$g_t(\mathbf{x}^*)=\phi_i(\mathbf{x}^*)$.
\end{proof}

\subsection{Scenario-aware Semantic Feedback IR}
\label{subsec:feedback-ir}

Raw failing tests omit the information the repair model needs:
\emph{which} scenario failed, \emph{why} the region matters (catch-all residual
or ordering precedence), and \emph{what} the oracle expected.
That information is organised into a structured semantic record---Semantic
Feedback IR---before any natural-language prompt is rendered.

The IR is designed so that scenario identity, guard text, observed/expected
outputs, and a typed violation reason travel together into the repair prompt,
rather than as an unstructured pass/fail dump.
Empirical comparisons of feedback surfaces (E6, E14, and related ablations)
are reported in Chapter~\ref{ch:results}; here we fix the representation and
the renderer contract.

\begin{equation}
  \mathcal{F}(i, c, t) = \langle \textit{type},\ s_i,\ g_i,\ \mathbf{x}_i^*,
  \ c(\mathbf{x}_i^*),\ g_t(\mathbf{x}_i^*),\ \textit{reason},\
  \textit{priority},\ \textit{fix} \rangle,
\end{equation}
with \textit{type} $\in \{\text{ordering}, \text{boundary}, \text{arithmetic},
\text{output}\}$, ranked \textit{priority}, and suggested \textit{fix}.
A template then renders the IR into a repair prompt that names the violated
scenario, states the guard, and states the expected postcondition.

\paragraph{Example (primary: catch-all).}
The buggy candidate of Listing~\ref{lst:others-bug} yields the JSON record
below (abbreviated).  The renderer projects the same object into the
natural-language quote of Section~\ref{sec:method:worked}:
\begin{lstlisting}[language={},basicstyle=\ttfamily\footnotesize,frame=single]
{
  "violation_type": "output",
  "scenario_index": 3,
  "constraint_text": "OTHERWISE (not Reject, not Pass)",
  "inputs": {"score": 3, "floor": 10},
  "observed": "Pass",
  "expected": "Review",
  "reason": "catch-all others returned wrong constant",
  "priority": "high",
  "fix": "return Review on residual band"
}
\end{lstlisting}
Serialisation follows \texttt{schemas/semantic\_feedback.schema.json}.
\texttt{FeedbackRenderer} projects the same IR onto the controlled E6 surfaces
(\texttt{test\_only}, \texttt{test\_expected}, full \texttt{semantic\_ir}) and
length-matched E14 surfaces (including \texttt{execution\_trace\_matched}),
separating failure structure from prompt text.
Default mode~M may render \texttt{execution\_trace\_matched} for repair;
E6 isolates \texttt{semantic\_ir} vs.\ \texttt{test\_only} as the primary C2
comparison (Chapter~\ref{ch:results}).

\paragraph{Example (secondary: ordering).}
On \code{classify\_signal}, an ordering failure on $(-3,-10)$ is recorded with
\texttt{violation\_type}: \texttt{ordering}, scenario index~1, observed
\texttt{Critical}, expected \texttt{Error}, and a fix hint that restores
$S_1$-first evaluation---the Semantic IR behind
Figure~\ref{fig:worked-example-trace}.

\subsection{Formal Conformance and Acceptance Predicate}
\label{subsec:conformance-acceptance}

Conformance answers a continuous question: on what fraction of witnesses does
the candidate match the oracle?
Acceptance answers a binary \emph{acceptance-safety} question: is the
candidate safe enough to ship under the conjunctive gate of
Figure~\ref{fig:accept-tree}---witnesses \emph{and} a structural screen?

\begin{equation}
  \label{eq:conformance}
  \mathrm{Conf}(c,t) = \frac{1}{|\mathcal{D}(t)|} \sum_{\mathbf{x} \in \mathcal{D}(t)}
  \mathbf{1}\bigl[f_c(\mathbf{x}) = g_t(\mathbf{x})\bigr].
\end{equation}
Let $\mathrm{Screen}(c,t)$ be the Stage~4 structural oracle
(Section~\ref{sec:method:screen}).
In the PoC reported here, $\mathrm{Screen}$ fails iff any smell in the
high-severity set $\mathcal{P}^H$ fires (Table~\ref{tab:patterns}); the
interface itself is not tied to that catalogue.
Release then requires both full witness agreement and a clear screen:
\begin{equation}
  \label{eq:accept}
  \mathrm{Accept}(c,t) \;\iff\;
  \mathrm{Conf}(c,t) = 1 \;\land\;
  \mathrm{Screen}(c,t) = \mathit{pass}.
\end{equation}
Equivalently, with the PoC encoding
$\mathrm{Screen}(c,t)=\mathit{pass} \iff
\sum_{p \in \mathcal{P}^H} \mathbf{1}[p(c,t)] = 0$.
On the secondary tax vignette, Equation~\eqref{eq:accept} rejects both a
foreign-high ordering bug exposed by $\Phi_2$ (Conf${<}1$) and an
unconditional \texttt{return "Low"} shortcut that might still look clean on
a lucky domestic suite (Screen fails via RF07).

\begin{proposition}[Finite-Sample FAR Bound under Independent Exposure]
\label{thm:far-bound}
\emph{Assumptions.}
Fix a faulty candidate and a task~$t$ with $|\mathcal{D}(t)|=n$ witnesses.
Assume (i)~each witness exposes the fault independently with probability at
least $\varepsilon>0$ (Definition~\ref{def:fault-exposure}), and
(ii)~witness draws are non-adversarial (i.i.d.\ relative to the fault).
Then
\[
  \mathrm{FAR} \;\leq\; (1-\varepsilon)^n.
\]
\end{proposition}
\begin{proof}
Write $E_i$ for the event that witness $i$ fails to expose the fault.
Independence and $\Pr[E_i]\le 1-\varepsilon$ give
\[
  \Pr[\text{all pass}\mid\text{faulty}]
  \;=\; \prod_{i=1}^{n}\Pr[E_i]
  \;\le\; (1-\varepsilon)^n.
\]
Hence $\mathrm{FAR}\le(1-\varepsilon)^n$.
\end{proof}

\begin{remark}[Scope relative to the deployed suite]
\label{rem:far-bound-scope}
Proposition~\ref{thm:far-bound} is a finite-sample envelope under the stated
i.i.d.\ exposure assumptions.
The deployed pipeline uses deterministic priority-encoded models
(Equation~\eqref{eq:priority}), not independent random draws; Proposition~\ref{thm:witness-soundness}
ensures each $\mathbf{x}\in\mathcal{D}(t)$ lies in a strict first-match region
under a faithful encoding, so $\varepsilon$ is bounded away from zero for the
mutation operators of Section~\ref{sec:exp:mutants}
(Corollary~\ref{cor:random-vs-smt}).
Acceptance under $\mathrm{Accept}(c,t)$ does not imply correctness on
$\mathcal{X}_t\setminus\mathcal{D}(t)$.
Witness cardinality also differs by stage: the repair loop uses up to
$N_{\max}$ cases (16 for mode~M); strict evaluation uses 64 cases; the
illustrative calculation in Section~\ref{subsec:far-validation} uses $n{=}9$
as an example only.
\end{remark}

\subsection{Structured Feedback and Repair Hypothesis}
\label{subsec:repair-monotonicity}

\begin{hypothesis}[Expected Improvement under Structured Feedback]
\label{hyp:structured-feedback}
Under fixed model and budget $K$, repair guided by Semantic Feedback IR
(which names the violated scenario, guard predicate, and violation type)
achieves higher expected conformance gain per iteration than repair guided
by plain test-pass/fail feedback.
\end{hypothesis}

\noindent
Chapter~\ref{ch:results} tests this hypothesis empirically
(Figures~\ref{fig:convergence} and~\ref{fig:feedback-variant}); a fully formal
stochastic model of LLM repair lies beyond current theory.

\subsection{Empirical FAR Envelope}
\label{subsec:far-validation}

With $\varepsilon = 0.08$ and illustrative $n = 9$ (not the per-repair cap
$N_{\max}{=}16$ or the strict-eval 64-case budget), Proposition~\ref{thm:far-bound} gives
$\mathrm{FAR} \leq (0.92)^9 \approx 0.473$ ($\approx 47\%$).
Chapter~\ref{ch:results} reports the operational E2 impl-screening rates
(mode~M FAR $5.0\%$, B2 $8.8\%$; PDR $95.0\%$ / $91.2\%$;
Table~\ref{tab:far-validation})---well inside that conservative envelope.

\begin{inlinetable}{Finite-sample FAR envelope vs.\ empirical E2 FAR
  (impl-screening; full protocol in Section~\ref{sec:results:rq3}).}
  \label{tab:far-validation}
  \footnotesize
  \begin{tabular}{lrrr}
    \toprule
    Mode & Empirical PDR & Empirical FAR & Bound $(1-\varepsilon)^n$ \\
    \midrule
    B2  & 91.2\% & 8.8\% & --- \\
    M   & 95.0\% & 5.0\% & $\approx 47\%$ \\
    \bottomrule
  \end{tabular}
\end{inlinetable}

\subsection{Prevention-Oriented Metrics}
\label{subsec:metrics}

For the faulty-mutant set $\mathcal{M}^F=\{m\in\mathcal{M}:\mathrm{Faulty}(m)=1\}$
(Definitions~\ref{def:pdr}--\ref{def:far}):
\begin{equation}
  \mathrm{PDR} = \frac{|\{m \in \mathcal{M}^F : \neg\mathrm{Accept}(c_m, t_m)\}|}{|\mathcal{M}^F|},
  \qquad
  \mathrm{FAR} = \frac{|\{m \in \mathcal{M}^F : \mathrm{Accept}(c_m, t_m)\}|}{|\mathcal{M}^F|}.
\end{equation}
PDR, FAR, strict formal conformance, and latency are reported jointly;
on a fixed faulty set, $\mathrm{PDR}+\mathrm{FAR}=1$.

\subsection{Pipeline Complexity Cost Model}
\label{subsec:complexity-cost}

With $n=|\mathcal{S}_t|$ scenarios and budget~$K$, witness generation costs
$\mathcal{O}(\sum_{i=1}^{n} T_{\text{SMT}}(|\Phi_i|))$ once per task; the
dominant per-run term is
$\mathcal{O}(K \cdot (T_{\text{LLM}} + n \cdot T_{\text{exec}}))$.
E3--E4 confirm sub-linear cost relative to conformance gain
(Corollary~\ref{cor:random-vs-smt}).

\begin{figure}[t]
  \centering
  \tikzinput{diagrams/tikz/prevention_protocol}
  \caption{Two-layer prevention evaluation: specification-confusion and
    implementation-screening mutant tracks aggregated to PDR and FAR.}
  \label{fig:prevention-protocol}
\end{figure}

\subsection{Full Algorithmic Procedure}
\label{sec:method:algorithm}

\begin{algorithm}[t]
\caption{HSP-Agile (\textsc{SgDP} instantiation)}
\label{alg:sgdp}
\begin{algorithmic}[1]
\Require task $t$, budget $K$
\Ensure accepted candidate or best-effort $\arg\max_k \mathrm{Conf}(c_k,t)$
\State $\mathcal{D}(t) \leftarrow \mathcal{W}(t)$
\State $\text{feedback}_0 \leftarrow \emptyset$; $\text{best} \leftarrow \mathtt{null}$
\For{$k = 1$ to $K$}
  \State $c_k \leftarrow \Fn{LLM}(t, \text{feedback}_{k-1})$
  \State $\mathrm{Conf}(c_k,t),\ \text{failures}_k \leftarrow$ eval $c_k$ on $\mathcal{D}(t)$ via $g_t$
  \State $\text{patterns}_k \leftarrow$ pattern guard on $c_k$
  \If{$\mathrm{Accept}(c_k, t)$} \Return $c_k$, accepted \EndIf
  \State $\mathcal{J}_k \leftarrow \textsc{ToJson}(\mathcal{F}(\text{failures}_k))$
        \Comment{schema: \texttt{semantic\_feedback.schema.json}}
  \State $\text{feedback}_k \leftarrow \rho(\mathcal{J}_k)$
        \Comment{\texttt{FeedbackRenderer} projection (E6 / E14 surfaces)}
  \State $\text{feedback}_k \mathrel{+}= \text{top-}5$ pattern violations from $\text{patterns}_k$
  \State update $\text{best}$ if $\mathrm{Conf}(c_k,t) > \mathrm{Conf}(\text{best},t)$
        \Comment{track $\arg\max$ Conf.\ over attempts}
  \State log $\langle k, \mathrm{Conf}(c_k,t), |\text{failures}_k|, |\text{patterns}_k|\rangle$
\EndFor
\Return $\text{best}$, not-accepted
        \Comment{best-effort: highest-Conf.\ candidate, not last}
\end{algorithmic}
\end{algorithm}

On budget exhaustion the implementation returns the best-effort candidate
$\arg\max_{k\le K}\mathrm{Conf}(c_k,t)$ among attempts (ties broken by formal
pass and fewer blocking pattern hits), not necessarily the last attempt $c_K$.
Each iteration also logs
$\langle k, \mathrm{Conf}(c_k,t), |\text{failures}_k|, |\text{patterns}_k|\rangle$
to \texttt{attempt\_history} for repair-dynamics analyses in
Chapter~\ref{ch:results}.

\begin{table*}[t]
\caption{Core Symbols and Definitions}
\label{tab:symbols}
\centering
\begin{tabular}{p{0.12\textwidth} p{0.82\textwidth}}
\toprule
Symbol & Meaning \\
\midrule
\(t\) & A benchmark task derived from Agile-SOFL specification or hard synthetic generator. \\
\((\mathcal{S}_t,\mathcal{X}_t,\mathcal{Y}_t)\) & Task tuple: ordered scenario sequence, input space, and output space for \(t\). \\
\(\mathcal{S}_t\) & Ordered FSF scenario set for task \(t\). \\
\(\mathcal{D}(t)\) & Concrete case set generated from \(\mathcal{S}_t\) by SMT-guided instantiation. \\
\(c, c_k\) & Candidate implementation generated by LLM; \(c_k\) denotes the attempt-\(k\) candidate. \\
\(f_c\) & Observable output function induced by candidate \(c\). \\
\(g_t\) & Specification-consistent oracle function for task \(t\). \\
\(K\) & Maximum synthesis-repair attempts allowed for one task run. \\
\(\tau\) & Strict success threshold for conformance (set to \(1\) in this study). \\
\(\mathrm{Conf}(c,t)\) & Strict formal conformance score, defined in Eq.~\ref{eq:conformance}; reported as \texttt{strict\_formal\_conformance} in artifacts. \\
\(\mathrm{StrictSuccess}(c,t)\) & Indicator that \(\mathrm{Conf}(c,t)\ge\tau\), Definition~\ref{def:strict-success}. \\
\(\mathrm{Accept}(c,t)\) & Hybrid acceptance predicate with formal+pattern constraints, Eq.~\ref{eq:accept}. \\
\(\mathcal{P}^{H}\) & High-severity requirement-fault pattern set. \\
\(\text{feedback}_k\) & Structured repair context after attempt \(k\) (counterexamples + pattern violations). \\
\(\mathcal{M}\) & Mutant population for prevention evaluation. \\
\(\mathrm{PDR}\) & Prevention detection rate over faulty mutants. \\
\(\mathrm{FAR}\) & False-accept / escape rate among faulty mutants. \\
\texttt{mutation\_kill\_rate} & Artifact metric for the proportion of mutants killed by a mode on the evaluated mutant population. \\
\bottomrule
\end{tabular}
\end{table*}

\begin{table*}[t]
\caption{Core Terminology}
\label{tab:glossary}
\centering
\begin{tabular}{p{0.28\textwidth} p{0.66\textwidth}}
\toprule
Term & Meaning \\
\midrule
\emph{Specification-conformance defect}
  & A fault where generated code behaves correctly on typical inputs yet
    violates the formal specification on precedence-critical or
    low-measure boundary regions that ordinary unit tests rarely reach. \\
\emph{First-match semantics}
  & Evaluation rule for ordered guard scenarios: the first scenario whose
    guard holds determines the output; later scenarios apply only to inputs
    not already claimed by earlier guards. \\
\emph{Witness}
  & A concrete input produced by SMT solving so that a designated scenario
    is the unique first match; used as an oracle test case. \\
\emph{Strict formal conformance}
  & Fraction of SMT witnesses on which the candidate matches the
    specification oracle ($\mathrm{Conf}(c,t)$). \\
\emph{Strict success / acceptance}
  & Release decision $\mathrm{Accept}(c,t)$: full witness conformance
    \emph{and} no high-severity pattern hits. Conformance can be high while
    acceptance remains low when residual witnesses or patterns still fail. \\
\emph{Semantic Feedback IR}
  & Typed JSON record of a conformance failure (scenario, guard, witness,
    observed vs.\ expected output, violation type) rendered into repair
    prompts. \\
\emph{SpecIR}
  & Notation-agnostic intermediate representation of ordered guarded cases;
    adapters map surface languages (FSF, Mini-Z, Mini-SM) into SpecIR. \\
\bottomrule
\end{tabular}
\end{table*}


\section{Design Principles}
\label{sec:method:principles}

Five principles~\citep{woodcock2009formalsurvey, chen2021codex}:
(P1)~oracle/SMT precision over heuristics;
(P2)~hard budget~$K$ with guaranteed termination;
(P3)~JSON Semantic Feedback IR before natural-language projection;
(P4)~conjunctive acceptance (equation~\eqref{eq:accept});
(P5)~modular stage flags for ablation (Chapter~\ref{ch:experiments}).

\section{Pipeline Overview}
\label{sec:method:overview}

The four stages of Algorithm~\ref{alg:sgdp} iterate until
$\mathrm{Accept}(c_k,t)=1$ or $k=K$.
Figure~\ref{fig:sgdp-overview} is the architectural map;
Figure~\ref{fig:pipeline-seq} is the message-level sequence of one run.

\begin{figure}[t]
  \centering
  \IfFileExists{diagrams/puml/rendered/pipeline_sequence.pdf}{%
    \pumlfig{diagrams/puml/rendered/pipeline_sequence}%
  }{%
    \IfFileExists{diagrams/puml/rendered/pipeline_sequence.png}{%
      \pumlfig{diagrams/puml/rendered/pipeline_sequence}%
    }{%
      \fbox{\parbox{0.75\textwidth}{\centering\vspace{0.8cm}
        \textit{[pipeline\_sequence.pdf pending]} \vspace{0.8cm}}}%
    }%
  }
  \caption{\textbf{Key observation:} one run is a message sequence until
    $\mathrm{Accept}$ or budget~$K$, not a single LLM call.
    \textbf{How to read:} synthesis $\to$ check $\to$ feedback $\to$ accept;
    the loop arc is Stage~3.}
  \label{fig:pipeline-seq}
\end{figure}

\section{Stage 1: LLM-Based Candidate Synthesis}
\label{sec:method:synthesis}

Stage~1 maps $(t, \mathit{feedback}_{k-1})$ to candidate $c_k$; it is the only
non-deterministic stage.

\subsection{Prompt Structure}
\label{subsec:method:prompt}

Prompt $\mathcal{P}(t,k)$ uses a fixed system message plus a user message with
$(\mathcal{S}_t, \mathit{sig}(t))$ and, when $k>1$, the rendered Feedback IR
\citep{openai2023gpt4, chen2021codex}.

\subsection{Model Configuration and Parsing}
\label{subsec:method:model}

Qwen2.5-27B-Instruct with $T_\text{gen}=0.7$ then $T_\text{repair}=0.0$
(Chapter~\ref{ch:experiments}); the parser extracts one fenced Python block
or sets $c_k=\bot$.
\begin{equation}
  c_k \;=\; \mathrm{LLM}\!\left(t,\;\mathit{feedback}_{k-1}\right).
  \label{eq:synthesis}
\end{equation}

\section{Stage 2: SMT-Driven Formal Conformance Checking}
\label{sec:method:checker}

Stage~2 realises $\mathcal{W}$ and $\mathrm{Conf}$
(equations~\eqref{eq:priority}--\eqref{eq:conformance}):
solve each $\Phi_i$, evaluate $g_t$ on the model, execute $c_k$ in a sandbox,
and collect mismatches as $\mathit{cexs}_k$~\citep{demoura2008z3, barrett2018smt}.
Algorithm~\ref{alg:checker} is the procedure; Listing~\ref{lst:z3-gen} shows
the Z3 encoding for one scenario.

\inputpython{listings/z3_gen.py}{Z3 test-case generation for scenario~$i$.}{lst:z3-gen}

\begin{algorithm}[t]
\caption{SMT-Driven Formal Conformance Checking}
\label{alg:checker}
\begin{algorithmic}[1]
\Require task $t$, candidate $c$
\Ensure $\mathrm{Conf}(c,t) \in [0,1]$, counterexample list $\mathit{cexs}$
\State $\mathcal{D}(t) \leftarrow \emptyset$; $\mathit{cexs} \leftarrow \emptyset$
\For{$i = 1$ \textbf{to} $n$}
  \State $\Phi_i \leftarrow \ProcCall{TranslateToZ3}{s_i,\, \{g_j\}_{j<i}}$
  \If{$\ProcCall{Z3-Solve}{\Phi_i} = \mathit{sat}$}
    \State $x_i \leftarrow \ProcCall{ExtractModel}{}$; $y_i^* \leftarrow \phi_i(x_i)$
    \State $\mathcal{D}(t) \leftarrow \mathcal{D}(t) \cup \{(x_i, y_i^*)\}$
  \EndIf
\EndFor
\For{\textbf{each} $(x, y^*) \in \mathcal{D}(t)$}
  \State $y^c \leftarrow f_c(x)$
  \If{$y^c \neq y^*$}
    \State $\mathit{cexs} \leftarrow \mathit{cexs} \cup
           \{(x, y^*, y^c, \ProcCall{ScenarioOf}{x})\}$
  \EndIf
\EndFor
\State \Return $(|\mathcal{D}(t)| - |\mathit{cexs}|)/|\mathcal{D}(t)|,\, \mathit{cexs})$
\end{algorithmic}
\end{algorithm}

\subsection{Why SMT Over Random Testing}
\label{subsec:method:smt-rationale}

Random sampling rarely hits multi-conjunct first-match regions or catch-all
residuals~\citep{claessen2000quickcheck, zhu1997specification}.
$\Phi_i$ places each witness in the exclusive region of scenario~$i$
(Proposition~\ref{thm:witness-soundness}), so oracle labels are unambiguous.

\section{Stage 3: Counterexample-Guided Repair}
\label{sec:method:repair}

Stage~3 builds the Semantic Feedback IR
(Section~\ref{subsec:feedback-ir}) and renders it via $\rho$;
Figure~\ref{fig:repair-feedback} shows the feedback arc for the primary
catch-all vignette (quote in Section~\ref{sec:method:worked}).
This is CEGIS with the LLM as synthesiser and the checker as
verifier~\citep{sun2024clover, solarlezama2008sketching}.

\begin{figure}[t]
  \centering
  \tikzinput{diagrams/tikz/repair_feedback}
  \caption{\textbf{Key observation (C2):} repair names the violated scenario
    and expected postcondition, not only raw I/O.
    \textbf{How to read:} failed attempt $\to$ Semantic Feedback IR $\to$ next
    synthesis.
    Controlled feedback comparisons are in Chapter~\ref{ch:results}.}
  \label{fig:repair-feedback}
\end{figure}

\section{Stage 4: Configurable Static Screen inside Dual-Gate Acceptance}
\label{sec:method:screen}
\label{sec:method:guard}

\subsection{Why Witnesses Alone Are Not Enough}
\label{subsec:method:why-screen}

Finite SMT witnesses answer a \emph{pointwise} question: does the candidate
match the oracle on forced first-match regions?
They do not answer a \emph{structural} question: did the model cheat the suite?
A degenerate candidate that always returns a constant success flag can survive
a thin suite that never exercises an overlap or a catch-all residual---the
LLM ``shortcut'' failure mode~\cite{dakhel2023copilot,hovemeyer2004findbugs}.
Acceptance safety requires a second conjunct after Conf.

Stage~4 is any oracle
\begin{equation}
  \mathrm{Screen} : (\mathit{code},\, t) \longrightarrow
  \{\mathit{pass},\mathit{fail}\}
\end{equation}
that emits severity-tagged hits into repair context.
Conjunctive Accept requires
$\mathrm{Conf}(c,t){=}1$ \emph{and}
$\mathrm{Screen}(c,t){=}\mathit{pass}$
(Figure~\ref{fig:accept-tree}; Equation~\eqref{eq:accept}).
On the ACL case (Section~\ref{sec:results:acl-case}), witnesses already refuse
the fail-open residual; Screen is the second latch when a sparse suite is
unlucky.

\subsection{Concrete Instantiation: AST Code-Smell Patterns (PoC)}
\label{subsec:method:smell-poc}

This study instantiates $\mathrm{Screen}$ as a catalogue of
requirements-level smells RF01--RF14~\citep{li2023iet_faultprevention}
(Appendix~\ref{app:patterns}).
Detectors operate primarily on the Python AST (control-flow and return shape),
with light surface checks where a smell is lexical.
Table~\ref{tab:patterns} lists the fourteen smells; nine high-severity ones
form the blocking set $\mathcal{P}^H$ in Equation~\eqref{eq:accept}, with RF07
(\emph{unconditional success}) as the canonical critical shortcut on
fail-closed contracts.
We treat this catalogue as one reproducible screen, not a completeness claim;
its operators draw on prior SOFL fault-prevention patterns, whereas
E2 also includes first-match and postcondition mutants that do not map
one-to-one onto these smells.
This separation prevents the prevention result from reducing to a detector
tested only on its own templates.

\begin{table}[H]
  \caption{PoC structural smell catalogue ($|\mathcal{P}| = 14$).
           Instance of the pluggable $\mathrm{Screen}$ interface.
           \textbf{critical} smells fail immediately; \textit{high} smells
           belong to the blocking set $\mathcal{P}^H$.}
  \label{tab:patterns}
  \centering
  \thesistable{%
    \small
    \setlength{\tabcolsep}{4pt}
    \begin{tabularx}{\linewidth}{@{}lXll@{}}
      \toprule
      ID   & Name                     & Category       & Severity \\
      \midrule
      RF01 & Missing precondition check  & requirement    & high \\
      RF02 & Wrong default branch        & FSF            & high \\
      RF03 & Hardcoded magic number      & implementation & medium \\
      RF04 & Missing output field        & interface      & high \\
      RF05 & Swapped comparison          & logic          & high \\
      RF06 & No guard for zero input     & boundary       & high \\
      RF07 & Unconditional success       & logic          & \textbf{critical} \\
      RF08 & Missing type coercion       & implementation & medium \\
      RF09 & Incomplete FSF coverage     & FSF            & high \\
      RF10 & External variable ignored   & SOFL           & medium \\
      RF11 & Off-by-one boundary         & boundary       & high \\
      RF12 & Duplicate assignment        & implementation & low \\
      RF13 & Empty handler               & logic          & high \\
      RF14 & Spec-comment mismatch       & requirement    & medium \\
      \bottomrule
    \end{tabularx}%
  }
\end{table}

\subsection{Configurable Screen Interface}
\label{subsec:method:pluggable-screen}

The evaluated mechanism is the second-gate role of a static screen inside
$\mathrm{Accept}$; RF01--RF14 are a reproducible PoC chosen for
alignment with prior SOFL fault-prevention patterns~\citep{li2023iet_faultprevention}.
Any component that consumes $(c,t)$ and emits severity-tagged hits---a project
smell pack, a CI static analyser
\cite{hovemeyer2004findbugs,ayewah2007findbugs}, or an LLM review
agent---can replace the PoC detectors without changing SpecIR witnesses,
Semantic Feedback IR, or Equation~\eqref{eq:accept}.

\subsection{Detection Method (PoC Wiring)}
\label{subsec:method:detection}

Each PoC smell $p \in \mathcal{P}$ has detector $\delta_p$; structural smells
use AST analysis, a minority use surface matching, and guard-comparison
smells cross-reference FSF guards with candidate conditionals.
$\mathrm{Hits}(c,t) = \{p : \delta_p(c,t)=1\}$; only $\mathcal{P}^H$ blocks
acceptance.
Replacing $\{\delta_p\}$ with another $\mathrm{Screen}$ implementation leaves
Equations~\eqref{eq:conformance}--\eqref{eq:accept} unchanged.

\section{Pipeline Stage Summary}
\label{sec:method:summary}

Table~\ref{tab:stage-summary} contracts the stages under Algorithm~\ref{alg:sgdp}.
The Tier-1 claims realised here are SpecIR semantic regions~(C1), Semantic
Feedback IR~(C2), and conjunctive acceptance~(C3).

\begin{inlinetable}{Summary of HSP-Agile pipeline stages.}
  \label{tab:stage-summary}
  \footnotesize
  \setlength{\tabcolsep}{3pt}
  \begin{tabularx}{\linewidth}{@{}clXXl@{}}
    \toprule
    Stage & Name          & Input               & Output                     & Property \\
    \midrule
    1     & Synthesis     & $t$, $\mathit{fb}$  & $c_k$                      & stochastic \\
    2     & Formal Check  & $c_k$, $t$          & $\mathrm{Conf}$, $\mathit{cexs}$ & sound (Prop.~\ref{thm:witness-soundness}) \\
    3     & Repair        & $\mathit{cexs}$, $\mathit{hits}$ & $\mathit{fb}_k$ & CEGIS bridge \\
    4     & Structural Screen & $c_k$, $t$          & $\mathit{hits}$ / veto     & pluggable ($\mathrm{Screen}$) \\
    ---   & Accept        & $\mathrm{Conf}$, $\mathrm{Screen}$ & $\{0,1\}$     & conjunctive (Prop.~\ref{thm:far-bound}) \\
    \bottomrule
  \end{tabularx}
\end{inlinetable}
